\chapter{Introducción y Objetivos}
\label{sec.ej4:introduccion}

El presente trabajo práctico aborda la simulación y análisis del proceso de muestreo y recuperación analógica de una
señal en banda base, así como la simulación de un sistema de codificación por modulación de impulsos codificados (PCM).

\section{Objetivos}
De acuerdo con la consigna, los objetivos fijados para la actividad son:
\begin{itemize}
    \item Realizar una simulación del proceso de muestreo y recuperación de una señal analógica mediante filtrado
          pasabajos.
    \item Validar la frecuencia de muestreo de acuerdo con el criterio de Nyquist.
    \item Mostrar y caracterizar el efecto de solapamiento espectral (aliasing) en el dominio temporal y frecuencial.
    \item Analizar el concepto de expansión del espectro derivado del muestreo.
    \item Determinar por qué el teorema resulta en una modulación de pulsos (PAM).
    \item Analizar el cuasi-muestreo a partir de muestrear con una portadora sin componente continua.
    \item Realizar una simulación de codificación PCM de una señal, verificando el comportamiento temporal en el bus
          paralelo de salida del convertidor y en la trama serie multiplexada.
\end{itemize}

\section{Esquemas de Implementación y Parámetros Fijados}
El desarrollo se estructuró a partir de dos esquemas de simulación en NI Multisim:
\begin{enumerate}
    \item \textbf{Muestreo y Recuperación de Señal:} La señal de banda base (punto A) se combina con la señal
          muestreadora (punto B) en un bloque muestreador para obtener la señal muestreada (punto C), la cual ingresa a
          una etapa de recuperación por filtrado pasabajos para obtener la señal analógica recuperada (punto D).
          Los valores de implementación fijados son:
          \begin{itemize}
              \item Señal de banda base: senoidal de $1\,\unit{\volt}$ de amplitud pico y frecuencia $10\,\unit{\kilo\hertz}$.
              \item Señal muestreadora: tren de pulsos con frecuencia nominal de $40\,\unit{\kilo\hertz}$.
          \end{itemize}
    \item \textbf{Sistema de Codificación PCM:} La señal analógica de entrada se muestrea y cuantiza mediante un
          conversor analógico-digital (ADC) de 8 bits, entregando la palabra binaria paralela en el punto de medición G.
          Dicha palabra se ingresa a un multiplexor digital para serializar los bits y obtener la trama binaria serie
          en el punto de medición H.
\end{enumerate}
